\documentclass[openany,fontset=fandol]{ctexbook}
\usepackage{amsmath,amssymb}
\usepackage[unicode,hidelinks]{hyperref}
\newcommand{\note}[1]{\footnote{译注：#1}}
\begin{document}
\title{代数几何基础第五卷预备笔记\\注13.4(c)：有限点集的邻域}
\author{据用户上传英译编辑稿译出}
\date{2026年11月17日}
\maketitle
\chapter*{范围与校订状态}
本批连续译出底本第53页注13.4(c)全文，包括编排建议、有限点集邻域的断言及其证明说明。下一处为同页13.4(d)：“By proceeding as explained in 7.1”。未合入原合订本。

本项两处余维数按同段的线性包维数及引理13.3校订，原写法和依据见脚注；七个问号代表的缺词照存，不猜补。闭点与$k$-有理点之间的过渡未在底本展开，译文不把二者等同。前批的8.8缺文、8.7(c)残缺及其他待核仍保留；本批不是对该推广或定理13.2的独立证明验证。
\chapter*{§13 初等态射与M. Artin定理（续）}
\noindent\textbf{注13.4（续）。}

\noindent c) （或可收入13.2的陈述。）沿用13.2的记号，设$\Phi$是$X$的一个有限子集，由$X$的光滑点组成，并假设$\Phi$包含在$X$的一个仿射开集中。则$\Phi$具有由多重初等开邻域组成的基本邻域系。显然，可以假设$\Phi$由闭点组成。

证明本质上相同，只须将13.3的表述略加改变如下：存在$P^r$中一个余维为$n$的线性子簇$C$，它不与$\Phi\cup Y$相交，并且对于每个$x_i\in\Phi(k)$，由$C$与$x_i$张成的余维为$n-1$的线性子簇$L_i$具有13.3的性质a)、b)、c)。\note{底本写“a linear subvariety $C$ of $P^r$ of codimension not meeting $\Phi\cup Y$”，在“codimension”之后漏去数值。同句规定$L_i=\langle C,x_i\rangle$的余维为$n-1$，且$x_i\notin C$，故$\dim L_i=\dim C+1$，所缺数值应为$n$；这也与第50页引理13.3之前的构造相符。底本“the linear subvariety of $L_i$”的多余“of”在正文删去。}\note{底本先说可假设$\Phi$由闭点组成，随后却写$x_i\in\Phi(k)$。闭点未必是$k$-有理点；此处照存记号，不擅将闭点改成有理点，也不增写未给出的基域扩张与下降论证。该过渡仍待核。}

为验证这一点，注意只须验证一般的$C$具有上述性质，因为对于这样的一个［原文缺词：???????］，每个$L_i$在经过$x_i$、余维为$n-1$的线性子簇$L$之中都是一般的，\note{此处底本写“among the $L$ of codimension $n$ passing through $x_i$”。然而同段的$L_i$及13.3中的$L_0$的余维均为$n-1$，故将这里的$n$订为$n-1$。对于固定的有理点$x_i$，每个经过$x_i$的这种$L$都可由其中不经过$x_i$的一个超平面$C$与$x_i$张成；这是与该余维数相容的线性代数说明，不是对有限点集情形中全部光滑性条件的独立证明。七个问号在底本中实际存在，其缺词不予恢复。}因此可以应用13.3的最初形式（或者至少应用我们已经证明的形式，即［法文标记］：每个经过$x$且充分一般的$L_0$都具有性质a)、b)、c)）。\note{底本末句保留“which was [Fr]”的编写痕迹；正文保留其法文标记，不臆造缺失的法语文字。这里“每个……充分一般的”是原文已证明形式的措辞，不应读成每一个经过$x$的线性子簇均满足这些性质。}
\end{document}
